\documentclass[12pt,a4paper]{article}
\usepackage[UTF8]{ctex}
\usepackage{amsmath,amssymb,amsthm}
\usepackage{geometry}
\usepackage{bm}
\usepackage{hyperref}
\usepackage{tikz}
\usetikzlibrary{arrows,arrows.meta,positioning,calc,shapes.geometric}

\geometry{margin=2.5cm}

\title{为什么Twist可以线性组合？}
\author{第12章抓取与操作补充说明}
\date{}

\newtheorem{theorem}{定理}
\newtheorem{definition}{定义}
\newtheorem{example}{例子}

\begin{document}

\maketitle

\section{问题}

在12.1.6章节中，我们经常看到这样的表达式：
\begin{equation}
V = k_1 V_1 + k_2 V_2 = (k_1 \omega_{1z} + k_2 \omega_{2z}, k_1 v_{1x} + k_2 v_{2x}, k_1 v_{1y} + k_2 v_{2y})
\end{equation}

\textbf{问题}：为什么twist可以这样线性组合？这有什么物理意义？

\section{答案：Twist空间是向量空间}

\subsection{向量空间的定义}

\textbf{关键事实}：Twist空间是一个\textbf{向量空间}（vector space），因此可以进行线性组合。

向量空间需要满足以下性质：
\begin{enumerate}
\item \textbf{加法封闭性}：如果$V_1$和$V_2$是twist，那么$V_1 + V_2$也是twist
\item \textbf{数乘封闭性}：如果$V$是twist，$k$是标量，那么$kV$也是twist
\item 满足向量空间的8条公理（结合律、交换律、分配律等）
\end{enumerate}

\subsection{为什么Twist空间是向量空间？}

\textbf{原因1}：Twist是速度的表示

\begin{itemize}
\item Twist表示刚体的\textbf{瞬时运动}
\item 包括角速度$\omega$和线速度$v$
\item 速度本身是向量，可以进行线性组合
\end{itemize}

\textbf{原因2}：速度的叠加原理

在物理学中，速度满足\textbf{叠加原理}：
\begin{itemize}
\item 如果物体同时有两个运动，总速度是两个速度的\textbf{向量和}
\item 这是经典力学的基本原理
\end{itemize}

\textbf{原因3}：数学结构

Twist $V = (\omega_z, v_x, v_y)$是$\mathbb{R}^3$中的向量：
\begin{itemize}
\item 每个分量都是实数
\item 可以进行标准的向量运算
\item 满足向量空间的所有性质
\end{itemize}

\section{线性组合的物理意义}

\subsection{基本理解}

\textbf{线性组合} $V = k_1 V_1 + k_2 V_2$的物理意义：

\begin{itemize}
\item $V_1$和$V_2$是两个可能的运动
\item $k_1$和$k_2$是"权重"系数
\item 组合后的运动$V$是这两个运动的"加权平均"
\end{itemize}

\subsection{具体例子}

\textbf{例子1}：$k_1 = 1, k_2 = 0$

\begin{itemize}
\item $V = 1 \cdot V_1 + 0 \cdot V_2 = V_1$
\item 就是运动$V_1$本身
\end{itemize}

\textbf{例子2}：$k_1 = 2, k_2 = 0$

\begin{itemize}
\item $V = 2 \cdot V_1$
\item 运动方向与$V_1$相同，但速度是$V_1$的2倍
\item 角速度：$\omega_z = 2 \omega_{1z}$
\item 线速度：$v = 2 v_1$
\end{itemize}

\textbf{例子3}：$k_1 = 0.5, k_2 = 0.5$

\begin{itemize}
\item $V = 0.5 V_1 + 0.5 V_2$
\item 这是两个运动的"平均"
\item 角速度：$\omega_z = 0.5 \omega_{1z} + 0.5 \omega_{2z}$
\item 线速度：$v = 0.5 v_1 + 0.5 v_2$
\end{itemize}

\section{为什么在接触约束中需要线性组合？}

\subsection{可行Twist集合是凸锥}

在接触约束中，可行twist集合通常是一个\textbf{凸锥}（convex cone）。

\textbf{凸锥的性质}：
\begin{itemize}
\item 如果$V_1$和$V_2$是可行twist
\item 那么它们的\textbf{非负线性组合} $k_1 V_1 + k_2 V_2$（$k_1, k_2 \geq 0$）也是可行twist
\item 这是凸锥的定义
\end{itemize}

\subsection{为什么是凸锥？}

\textbf{原因}：接触约束是\textbf{线性不等式}：
\begin{equation}
F^T V \geq 0
\end{equation}

如果$V_1$和$V_2$都满足约束：
\begin{align}
F^T V_1 &\geq 0 \\
F^T V_2 &\geq 0
\end{align}

那么对于$k_1, k_2 \geq 0$：
\begin{align}
F^T (k_1 V_1 + k_2 V_2) &= k_1 F^T V_1 + k_2 F^T V_2 \\
&\geq k_1 \cdot 0 + k_2 \cdot 0 = 0
\end{align}

所以$k_1 V_1 + k_2 V_2$也满足约束！

\section{详细数学推导}

\subsection{Twist的表示}

平面twist：
\begin{equation}
V = \begin{bmatrix} \omega_z \\ v_x \\ v_y \end{bmatrix} \in \mathbb{R}^3
\end{equation}

这是$\mathbb{R}^3$中的标准向量。

\subsection{线性组合的定义}

对于两个twist：
\begin{align}
V_1 &= \begin{bmatrix} \omega_{1z} \\ v_{1x} \\ v_{1y} \end{bmatrix} \\
V_2 &= \begin{bmatrix} \omega_{2z} \\ v_{2x} \\ v_{2y} \end{bmatrix}
\end{align}

它们的线性组合：
\begin{align}
V &= k_1 V_1 + k_2 V_2 \\
&= k_1 \begin{bmatrix} \omega_{1z} \\ v_{1x} \\ v_{1y} \end{bmatrix} + k_2 \begin{bmatrix} \omega_{2z} \\ v_{2x} \\ v_{2y} \end{bmatrix} \\
&= \begin{bmatrix} k_1 \omega_{1z} + k_2 \omega_{2z} \\ k_1 v_{1x} + k_2 v_{2x} \\ k_1 v_{1y} + k_2 v_{2y} \end{bmatrix}
\end{align}

这是标准的向量运算，完全合法！

\subsection{为什么这样组合有意义？}

\textbf{物理原因}：速度的叠加

\begin{itemize}
\item 如果物体可以以速度$V_1$运动
\item 也可以以速度$V_2$运动
\item 那么它也可以以"混合"速度$k_1 V_1 + k_2 V_2$运动
\item 这是速度叠加原理的直接应用
\end{itemize}

\textbf{数学原因}：向量空间结构

\begin{itemize}
\item $\mathbb{R}^3$是向量空间
\item 任何向量空间中的元素都可以线性组合
\item Twist是$\mathbb{R}^3$中的向量
\item 因此可以线性组合
\end{itemize}

\section{正线性组合的特殊意义}

\subsection{为什么需要$k_1, k_2 \geq 0$？}

在接触约束中，我们通常考虑\textbf{正线性组合}（非负线性组合）：
\begin{equation}
V = k_1 V_1 + k_2 V_2, \quad k_1, k_2 \geq 0
\end{equation}

\textbf{原因1}：可行集合是凸锥

\begin{itemize}
\item 可行twist集合是凸锥
\item 凸锥由非负线性组合生成
\item 所以只需要考虑$k_1, k_2 \geq 0$
\end{itemize}

\textbf{原因2}：物理意义

\begin{itemize}
\item $k_1, k_2 \geq 0$意味着"正向"组合
\item 不会"抵消"运动
\item 保持运动的"方向性"
\end{itemize}

\textbf{原因3}：约束的保持

\begin{itemize}
\item 如果$V_1$和$V_2$都满足约束$F^T V \geq 0$
\item 那么$k_1 V_1 + k_2 V_2$（$k_1, k_2 \geq 0$）也满足约束
\item 这是凸锥的性质
\end{itemize}

\section{具体应用：CoR的线性组合}

\subsection{问题回顾}

给定两个twist $V_1$和$V_2$，它们的线性组合$V = k_1 V_1 + k_2 V_2$对应的CoR在哪里？

\subsection{推导过程}

\textbf{步骤1}：计算组合后的twist

\begin{equation}
V = (k_1 \omega_{1z} + k_2 \omega_{2z}, k_1 v_{1x} + k_2 v_{2x}, k_1 v_{1y} + k_2 v_{2y})
\end{equation}

\textbf{步骤2}：计算CoR位置

如果$\omega_z = k_1 \omega_{1z} + k_2 \omega_{2z} \neq 0$：
\begin{align}
\text{CoR}(V) &= \left(-\frac{v_y}{\omega_z}, \frac{v_x}{\omega_z}\right) \\
&= \left(-\frac{k_1 v_{1y} + k_2 v_{2y}}{k_1 \omega_{1z} + k_2 \omega_{2z}}, \frac{k_1 v_{1x} + k_2 v_{2x}}{k_1 \omega_{1z} + k_2 \omega_{2z}}\right)
\end{align}

\textbf{关键观察}：这个CoR位置依赖于$k_1$和$k_2$的比值，而不是绝对值。

\subsection{为什么CoR在直线上？}

\textbf{数学证明}：

设$t = \frac{k_1}{k_1 + k_2}$（当$k_1 + k_2 \neq 0$时）

则：
\begin{align}
\text{CoR}(V) &= \left(-\frac{t v_{1y} + (1-t) v_{2y}}{t \omega_{1z} + (1-t) \omega_{2z}}, \frac{t v_{1x} + (1-t) v_{2x}}{t \omega_{1z} + (1-t) \omega_{2z}}\right)
\end{align}

当$t$从0变化到1时，CoR在通过$\text{CoR}(V_1)$和$\text{CoR}(V_2)$的直线上移动。

\section{常见疑问解答}

\subsection{疑问1：线性组合的物理实现}

\textbf{问题}：物体真的可以同时进行两个运动吗？

\textbf{回答}：
\begin{itemize}
\item 不是"同时"进行两个运动
\item 而是"可以选择"进行组合后的运动
\item 线性组合表示的是\textbf{可能的运动集合}
\item 物体在某个时刻只能有一个运动，但可以选择组合后的运动
\end{itemize}

\subsection{疑问2：为什么可以"混合"角速度和线速度？}

\textbf{问题}：角速度和线速度是不同的物理量，为什么可以混合？

\textbf{回答}：
\begin{itemize}
\item 虽然物理意义不同，但数学上都是实数
\item Twist $(\omega_z, v_x, v_y)$是$\mathbb{R}^3$中的向量
\item 向量空间的元素可以线性组合
\item 组合后的结果仍然是有效的twist
\end{itemize}

\subsection{疑问3：线性组合是否总是有物理意义？}

\textbf{问题}：任何线性组合都对应实际可能的运动吗？

\textbf{回答}：
\begin{itemize}
\item \textbf{数学上}：任何线性组合都是有效的twist
\item \textbf{物理上}：需要满足约束条件
\item 在接触约束中，只有满足$F^T V \geq 0$的twist是可行的
\item 线性组合帮助我们找到所有可行的twist
\end{itemize}

\section{总结}

\subsection{关键点}

\begin{enumerate}
\item \textbf{Twist空间是向量空间}：可以进行线性组合

\item \textbf{线性组合的物理意义}：
  \begin{itemize}
  \item 速度的叠加
  \item 可能的运动集合
  \item 凸锥的生成
  \end{itemize}

\item \textbf{在接触约束中的应用}：
  \begin{itemize}
  \item 可行twist集合是凸锥
  \item 由可行twist的非负线性组合生成
  \item 用于分析CoR的位置
  \end{itemize}

\item \textbf{数学基础}：
  \begin{itemize}
  \item 向量空间理论
  \item 凸分析
  \item 线性代数
  \end{itemize}
\end{enumerate}

\subsection{公式总结}

\textbf{线性组合}：
\begin{equation}
V = k_1 V_1 + k_2 V_2 = (k_1 \omega_{1z} + k_2 \omega_{2z}, k_1 v_{1x} + k_2 v_{2x}, k_1 v_{1y} + k_2 v_{2y})
\end{equation}

\textbf{正线性组合}（$k_1, k_2 \geq 0$）：
\begin{equation}
V = \text{pos}(\{V_1, V_2\}) = \{k_1 V_1 + k_2 V_2 | k_1, k_2 \geq 0\}
\end{equation}

\textbf{CoR位置}（当$\omega_z \neq 0$时）：
\begin{equation}
\text{CoR}(V) = \left(-\frac{v_y}{\omega_z}, \frac{v_x}{\omega_z}\right)
\end{equation}

\section{重要补充：坐标系必须相同}

\subsection{你的观察}

\textbf{关键问题}：如果要进行线性组合 $V = k_1 V_1 + k_2 V_2$，$V_1$和$V_2$必须在\textbf{同一个坐标系}中表示！

\subsection{为什么必须相同？}

\textbf{原因1}：向量运算的前提

\begin{itemize}
\item 线性组合是\textbf{分量对应相加}
\item 如果坐标系不同，分量没有对应关系
\item 例如：$V_1$的$x$分量和$V_2$的$x$分量可能指向不同方向
\item 直接相加没有意义
\end{itemize}

\textbf{原因2}：物理意义的一致性

\begin{itemize}
\item 如果$V_1$和$V_2$在不同坐标系中
\item 它们的"方向"含义不同
\item 组合后的结果无法解释
\end{itemize}

\textbf{原因3}：数学运算的合法性

\begin{itemize}
\item 向量空间的线性组合要求元素在\textbf{同一个空间}中
\item 不同坐标系中的twist属于不同的向量空间
\item 必须先转换到同一坐标系
\end{itemize}

\subsection{在接触约束中的应用}

在12.1.3和12.1.4章节中：

\begin{itemize}
\item 所有twist都在\textbf{空间坐标系}（space frame）$\{s\}$中表示
\item 这是第12章的标准约定
\item 因此可以安全地进行线性组合
\end{itemize}

\textbf{注意}：在第12章中，所有twist和wrench都在空间坐标系中表示（见12.1.2节的脚注）。

\subsection{如果坐标系不同怎么办？}

如果$V_1$在坐标系$\{a\}$中，$V_2$在坐标系$\{b\}$中：

\begin{enumerate}
\item \textbf{步骤1}：将$V_1$转换到$\{b\}$，或
\item \textbf{步骤2}：将$V_2$转换到$\{a\}$，或
\item \textbf{步骤3}：将两者都转换到第三个坐标系$\{c\}$
\end{enumerate}

然后才能进行线性组合。

\textbf{转换公式}（回顾第3章）：

如果$T_{ab}$是从$\{a\}$到$\{b\}$的变换，则：
\begin{equation}
V_b = [\text{Ad}_{T_{ab}}] V_a
\end{equation}

其中$[\text{Ad}_{T_{ab}}]$是伴随变换矩阵。

\subsection{具体例子}

\textbf{例子}：两个twist在不同坐标系中

\begin{itemize}
\item $V_1$在坐标系$\{1\}$中：$V_1 = (1, 2, 0)$
\item $V_2$在坐标系$\{2\}$中：$V_2 = (1, 0, 2)$
\item \textbf{错误}：直接计算 $V = V_1 + V_2$（没有意义！）
\item \textbf{正确}：先将$V_1$转换到$\{2\}$，然后相加
\end{itemize}

\subsection{在CoR计算中的意义}

当我们计算两个twist的线性组合对应的CoR时：

\begin{itemize}
\item 两个twist必须在\textbf{同一个坐标系}中
\item CoR的位置也在\textbf{同一个坐标系}中
\item 这样才能正确理解CoR的几何位置
\end{itemize}

\subsection{总结}

\begin{enumerate}
\item \textbf{线性组合要求}：所有twist必须在同一坐标系中

\item \textbf{第12章的约定}：所有twist在空间坐标系$\{s\}$中表示

\item \textbf{如果坐标系不同}：必须先进行坐标变换

\item \textbf{物理意义}：只有在同一坐标系中，线性组合才有明确的物理意义
\end{enumerate}

\textbf{你的观察非常正确！}这是理解twist线性组合的重要前提。

\section{补充：CoR位置的凸组合表示（即使符号不同）}

\subsection{你的问题}

即使角速度符号不同（$\omega_{1z}$和$\omega_{2z}$符号相反），CoR位置是否也能写成：
\begin{equation}
\text{CoR}(V) = \frac{k_1}{k_1 + k_2} \text{CoR}_1 + \frac{k_2}{k_1 + k_2} \text{CoR}_2
\label{eq:cor_convex}
\end{equation}

\subsection{答案：可以，但需要小心}

\textbf{简短答案}：数学上可以写成这种形式，但需要理解其几何意义。

\subsection{数学推导}

\textbf{步骤1}：CoR的基本公式

对于twist $V = (\omega_z, v_x, v_y)$，如果$\omega_z \neq 0$：
\begin{equation}
\text{CoR}(V) = \left(-\frac{v_y}{\omega_z}, \frac{v_x}{\omega_z}\right)
\end{equation}

\textbf{步骤2}：线性组合的CoR

对于$V = k_1 V_1 + k_2 V_2$：
\begin{align}
V &= (k_1 \omega_{1z} + k_2 \omega_{2z}, k_1 v_{1x} + k_2 v_{2x}, k_1 v_{1y} + k_2 v_{2y}) \\
\omega_z &= k_1 \omega_{1z} + k_2 \omega_{2z}
\end{align}

如果$\omega_z \neq 0$：
\begin{align}
\text{CoR}(V) &= \left(-\frac{k_1 v_{1y} + k_2 v_{2y}}{k_1 \omega_{1z} + k_2 \omega_{2z}}, \frac{k_1 v_{1x} + k_2 v_{2x}}{k_1 \omega_{1z} + k_2 \omega_{2z}}\right)
\end{align}

\textbf{步骤3}：引入参数$t$

设$t = \frac{k_1}{k_1 + k_2}$（当$k_1 + k_2 \neq 0$时）

则：
\begin{align}
k_1 &= t(k_1 + k_2) \\
k_2 &= (1-t)(k_1 + k_2)
\end{align}

\textbf{步骤4}：重写CoR公式

\begin{align}
\text{CoR}(V) &= \left(-\frac{t(k_1+k_2) v_{1y} + (1-t)(k_1+k_2) v_{2y}}{t(k_1+k_2) \omega_{1z} + (1-t)(k_1+k_2) \omega_{2z}}, \frac{t(k_1+k_2) v_{1x} + (1-t)(k_1+k_2) v_{2x}}{t(k_1+k_2) \omega_{1z} + (1-t)(k_1+k_2) \omega_{2z}}\right) \\
&= \left(-\frac{t v_{1y} + (1-t) v_{2y}}{t \omega_{1z} + (1-t) \omega_{2z}}, \frac{t v_{1x} + (1-t) v_{2x}}{t \omega_{1z} + (1-t) \omega_{2z}}\right)
\end{align}

\textbf{关键观察}：这个公式\textbf{不能}直接写成$\text{CoR}_1$和$\text{CoR}_2$的凸组合，因为分母是$\omega_z$，不是$\omega_{1z}$或$\omega_{2z}$。

\subsection{特殊情况：当$\omega_{1z}$和$\omega_{2z}$同号时}

如果$\omega_{1z}$和$\omega_{2z}$同号（都为正或都为负），且$k_1, k_2 \geq 0$：

\begin{itemize}
\item $\omega_z = t \omega_{1z} + (1-t) \omega_{2z}$与$\omega_{1z}$和$\omega_{2z}$同号
\item 可以证明（需要一些代数运算）：
  \begin{equation}
  \text{CoR}(V) = \frac{t \omega_{1z}}{t \omega_{1z} + (1-t) \omega_{2z}} \text{CoR}_1 + \frac{(1-t) \omega_{2z}}{t \omega_{1z} + (1-t) \omega_{2z}} \text{CoR}_2
  \end{equation}
\item 这是\textbf{加权平均}，权重与角速度成正比
\item 当$\omega_{1z} = \omega_{2z}$时，简化为标准的凸组合
\end{itemize}

\subsection{一般情况：符号不同时}

当$\omega_{1z}$和$\omega_{2z}$符号不同时：

\textbf{问题}：$\omega_z = t \omega_{1z} + (1-t) \omega_{2z}$的符号会改变！

\begin{itemize}
\item 当$t$接近1时：$\omega_z \approx \omega_{1z}$（与$\omega_{1z}$同号）
\item 当$t$接近0时：$\omega_z \approx \omega_{2z}$（与$\omega_{2z}$同号）
\item 当$t = t_0 = \frac{-\omega_{2z}}{\omega_{1z} - \omega_{2z}}$时：$\omega_z = 0$（纯平移）
\end{itemize}

\textbf{CoR位置的变化}：

\begin{itemize}
\item 当$t > t_0$时：$\omega_z$与$\omega_{1z}$同号，CoR在从CoR$_1$出发的射线上
\item 当$t < t_0$时：$\omega_z$与$\omega_{2z}$同号，CoR在从CoR$_2$出发的射线上
\item 当$t = t_0$时：$\omega_z = 0$，CoR在无穷远处（平移）
\end{itemize}

\textbf{几何解释}：

虽然不能写成简单的凸组合，但CoR仍然在通过CoR$_1$和CoR$_2$的\textbf{整条直线}上，只是：
\begin{itemize}
\item 中间有一个"跳跃"（对应$\omega_z = 0$的点）
\item 这个点将直线分成两部分
\item 每部分对应不同的旋转方向（'+'或'-'）
\end{itemize}

\subsection{你提到的公式的适用条件}

你提到的公式：
\begin{equation}
\text{CoR}(V) = \frac{k_1}{k_1 + k_2} \text{CoR}_1 + \frac{k_2}{k_1 + k_2} \text{CoR}_2
\end{equation}

\textbf{适用条件}：
\begin{enumerate}
\item $\omega_{1z}$和$\omega_{2z}$同号（都为正或都为负）
\item 或者，更一般地，当$\omega_z = k_1 \omega_{1z} + k_2 \omega_{2z}$与$\omega_{1z}$和$\omega_{2z}$同号时
\item 这是\textbf{加权平均}，权重是$\frac{k_1}{k_1 + k_2}$和$\frac{k_2}{k_1 + k_2}$
\end{enumerate}

\textbf{不适用的情况}：
\begin{itemize}
\item 当$\omega_{1z}$和$\omega_{2z}$符号不同，且$\omega_z$的符号会改变时
\item 此时CoR不在两个CoR之间的线段上，而是在两条射线上
\end{itemize}

\subsection{总结}

\begin{enumerate}
\item \textbf{数学上}：CoR位置可以写成某种形式的组合，但不总是简单的凸组合

\item \textbf{相同符号}：可以写成加权平均（权重与角速度成正比）

\item \textbf{不同符号}：CoR在整条直线上，但被$\omega_z = 0$的点分成两部分

\item \textbf{几何意义}：CoR总是在通过CoR$_1$和CoR$_2$的直线上，这是线性组合的直接结果
\end{enumerate}

\textbf{你的公式在特定条件下是正确的！}特别是当角速度同号时。

\section{什么是凸组合（Convex Combination）？}

\subsection{定义}

\begin{definition}[凸组合]
给定$n$个向量（或点）$x_1, x_2, \ldots, x_n$，它们的\textbf{凸组合}定义为：
\begin{equation}
x = \sum_{i=1}^n \lambda_i x_i
\end{equation}
其中系数$\lambda_i$满足：
\begin{enumerate}
\item \textbf{非负性}：$\lambda_i \geq 0$ 对所有$i$
\item \textbf{归一性}：$\sum_{i=1}^n \lambda_i = 1$
\end{enumerate}
\end{definition}

\subsection{关键特征}

凸组合的两个关键条件：

\begin{enumerate}
\item \textbf{系数非负}：$\lambda_i \geq 0$
  \begin{itemize}
  \item 不能有负系数
  \item 这确保组合结果在"内部"，不会"超出"原始点的范围
  \end{itemize}

\item \textbf{系数和为1}：$\sum \lambda_i = 1$
  \begin{itemize}
  \item 这是"加权平均"
  \item 确保组合结果是原始点的"混合"，而不是缩放
  \end{itemize}
\end{enumerate}

\subsection{最简单的例子：两个点的凸组合}

对于两个点$x_1$和$x_2$，凸组合为：
\begin{equation}
x = \lambda x_1 + (1-\lambda) x_2, \quad 0 \leq \lambda \leq 1
\end{equation}

\textbf{几何意义}：
\begin{itemize}
\item 当$\lambda = 0$时：$x = x_2$
\item 当$\lambda = 1$时：$x = x_1$
\item 当$\lambda = 0.5$时：$x = 0.5 x_1 + 0.5 x_2$（中点）
\item 当$\lambda$从0变化到1时：$x$在从$x_2$到$x_1$的\textbf{线段}上移动
\end{itemize}

\textbf{图示}：

\begin{center}
\begin{tikzpicture}[scale=2]
% 两个点
\filldraw[red] (0,0) circle (2pt) node[below left] {$x_1$};
\filldraw[blue] (3,1) circle (2pt) node[below right] {$x_2$};

% 连接线段
\draw[thick,green] (0,0) -- (3,1);

% 标记几个凸组合的点
\filldraw[purple] (0.75,0.25) circle (1.5pt) node[above] {$\lambda=0.25$};
\filldraw[purple] (1.5,0.5) circle (1.5pt) node[above] {$\lambda=0.5$};
\filldraw[purple] (2.25,0.75) circle (1.5pt) node[above] {$\lambda=0.75$};

\node[below] at (1.5,-0.3) {凸组合：连接两个点的线段};
\end{tikzpicture}
\end{center}

\subsection{三个点的凸组合}

对于三个点$x_1, x_2, x_3$，凸组合为：
\begin{equation}
x = \lambda_1 x_1 + \lambda_2 x_2 + \lambda_3 x_3
\end{equation}
其中：$\lambda_1, \lambda_2, \lambda_3 \geq 0$且$\lambda_1 + \lambda_2 + \lambda_3 = 1$

\textbf{几何意义}：
\begin{itemize}
\item 凸组合的结果在三个点形成的\textbf{三角形内部}（包括边界）
\item 这是三个点的\textbf{凸包}（convex hull）
\end{itemize}

\begin{center}
\begin{tikzpicture}[scale=2]
% 三个点
\filldraw[red] (0,0) circle (2pt) node[below left] {$x_1$};
\filldraw[blue] (3,0) circle (2pt) node[below right] {$x_2$};
\filldraw[green] (1.5,2) circle (2pt) node[above] {$x_3$};

% 三角形
\fill[purple!20] (0,0) -- (3,0) -- (1.5,2) -- cycle;
\draw[thick,green] (0,0) -- (3,0) -- (1.5,2) -- cycle;

\node[purple] at (1.5,0.7) {凸组合区域\\（三角形）};
\end{tikzpicture}
\end{center}

\subsection{凸组合 vs 线性组合 vs 正线性组合}

\begin{center}
\begin{tabular}{|c|c|c|c|}
\hline
\textbf{类型} & \textbf{系数条件} & \textbf{几何意义} & \textbf{例子} \\
\hline
线性组合 & $k_i \in \mathbb{R}$ & 整个空间 & 任何方向 \\
\hline
正线性组合 & $k_i \geq 0$ & 凸锥 & 从原点出发的射线 \\
\hline
凸组合 & $k_i \geq 0$且$\sum k_i = 1$ & 凸包（有界） & 线段、三角形等 \\
\hline
\end{tabular}
\end{center}

\textbf{关系}：
\begin{equation}
\text{凸组合} \subseteq \text{正线性组合} \subseteq \text{线性组合}
\end{equation}

\subsection{在CoR中的应用}

\subsubsection{相同符号角速度的情况}

当$\omega_{1z}$和$\omega_{2z}$同号时，CoR位置可以写成：
\begin{equation}
\text{CoR}(V) = \frac{k_1}{k_1 + k_2} \text{CoR}_1 + \frac{k_2}{k_1 + k_2} \text{CoR}_2
\end{equation}

\textbf{这是凸组合吗？}

\begin{itemize}
\item 系数：$\lambda_1 = \frac{k_1}{k_1 + k_2}$，$\lambda_2 = \frac{k_2}{k_1 + k_2}$
\item 检查条件1：$\lambda_1, \lambda_2 \geq 0$？是的（因为$k_1, k_2 \geq 0$）
\item 检查条件2：$\lambda_1 + \lambda_2 = 1$？是的
\item \textbf{结论}：这是凸组合！
\end{itemize}

\textbf{几何意义}：
\begin{itemize}
\item CoR在连接CoR$_1$和CoR$_2$的\textbf{线段}上
\item 这是凸组合的典型特征
\end{itemize}

\subsubsection{不同符号角速度的情况}

当$\omega_{1z}$和$\omega_{2z}$符号不同时：
\begin{itemize}
\item 不能写成简单的凸组合
\item CoR不在两个CoR之间的线段上
\item 而是在两条射线上（分别从两个CoR出发）
\end{itemize}

\textbf{为什么不是凸组合？}

\begin{itemize}
\item 虽然可以写成某种形式，但系数不满足凸组合的条件
\item 或者，更准确地说，CoR位置不在"线段"上，而是在"直线"上
\item 凸组合的结果总是在有界区域内，但这里CoR可以延伸到无穷远
\end{itemize}

\subsection{凸组合的重要性}

\subsubsection{在接触约束中}

\begin{itemize}
\item \textbf{可行twist集合}：通常是凸锥（正线性组合）
\item \textbf{形状闭合}：要求接触wrench的凸包包含原点
\item \textbf{CoR区域}：多个接触的可行CoR区域是凸的
\end{itemize}

\subsubsection{在优化中}

\begin{itemize}
\item 凸组合保持凸性
\item 如果约束是凸的，可行区域也是凸的
\item 这简化了优化问题的求解
\end{itemize}

\subsection{总结}

\begin{enumerate}
\item \textbf{凸组合}：系数非负且和为1的线性组合

\item \textbf{几何意义}：结果在原始点的凸包内（有界区域）

\item \textbf{两个点的凸组合}：连接两个点的线段

\item \textbf{在CoR中}：
  \begin{itemize}
  \item 相同符号：CoR是凸组合，在线段上
  \item 不同符号：CoR不是凸组合，在射线上
  \end{itemize}

\item \textbf{重要性}：凸组合是理解可行区域几何结构的关键
\end{enumerate}

\end{document}

